\chapter{Systemarchitektur und Technologie}
\label{chap:arch}

\section{DApp-Kernkonzepte und ihre Umsetzung im Projekt}

Die Lehrveranstaltung verlangt typischerweise die sichere Anwendung
grundlegender DApp-Bausteine: eine öffentliche Blockchain als Vertrauensanker,
in Solidity geschriebene Smart Contracts, ein Browser-Frontend und eine Wallet,
die Transaktionen signiert. Im Merkblatt werden beispielhaft Web3.js und MetaMask
genannt; funktional äquivalent setzt dieses Projekt auf \textbf{ethers.js in der
Version 6} und ebenfalls MetaMask als injizierten Ethereum-Provider
(\texttt{window.ethereum}).

Das Frontend erzeugt für jeden benötigten Vertrag eine \texttt{Contract}-Instanz
mit ABI und Deployment-Adresse. Lesezugriffe nutzen den Provider; schreibende
Aufrufe laufen über einen \texttt{Signer}, der von der verbundenen Wallet
kommt. Damit ist die kryptographische Autorisierung klar von der bloßen
Datenabfrage getrennt: Nur wenn Nutzerinnen eine Transaktion bestätigen, ändert
sich der On-Chain-Zustand.

\section{Schichtenmodell und Datenflüsse}

\begin{sloppypar}
Die Architektur lässt sich als dreischichtiges Modell beschreiben, das von der
Präsentation bis zur persistenten Chain reicht. In komprimierter Lesart ergibt
sich folgender Datenfluss: Die Nuxt-Oberfläche ruft über \texttt{ethers.js}~v6
die MetaMask-Provider-Schnittstelle auf; MetaMask signiert Transaktionen und
sendet sie an einen JSON-RPC-Endpunkt (lokal Hardhat, öffentlich z.\,B.\
Sepolia). Lesende Aufrufe treffen dieselben Verträge \texttt{VoteToken} und
\texttt{PredictionMarket} ohne Wallet-Bestätigung, sofern die Anwendung einen
reinen Provider-Pfad nutzt---im vorliegenden Frontend dominiert jedoch der
Signer-Pfad für Einheitlichkeit.
\end{sloppypar}

\paragraph{Präsentationsschicht.}
Das Verzeichnis \path{frontend/} enthält die Nuxt-Anwendung. Die zentrale
Oberfläche liegt in \path{frontend/app/app.vue}; wiederkehrende Logik ist in
Composables ausgelagert, insbesondere \path{useWallet.ts} für Verbindung, Kette
und Signer sowie \path{useContracts.ts} für Vertragsinstanzen, Aggregationslese-
routinen und Transaktionshüllen.

\paragraph{On-Chain-Logik.}
Die Verträge \path{contracts/VoteToken.sol} und
\path{contracts/PredictionMarket.sol} bilden die fachliche Wahrheit. Sie werden
mit Hardhat kompiliert, getestet und über \path{scripts/deploy.ts} ausgerollt.
Die Trennung zwischen Token und Markt folgt dem Prinzip der getrennten
Zuständigkeiten: Der Token kapselt Ausgabe und Übertragung nach ERC-20; der
Marktvertrag kapselt Geschäftsregeln und hält aggregierte Pools.

\paragraph{Build- und Integrationspfad.}
Nach dem Deployment schreibt \path{scripts/deploy.ts} die Adressen nach
\path{deployments/<netzwerk>.json}. Das Skript \path{scripts/export-frontend.ts}
überträgt ABIs und Adressen in \path{frontend/app/contracts/}, sodass die
Oberfläche ohne manuelles Kopieren mit der jeweiligen Deployment-Generation
übereinstimmt. Dieser Schritt ist leicht zu vergessen, aber für reproduzierbare
Demos unverzichtbar.

\section{Abwesenheit eines zentralen Anwendungsservers für die Geschäftslogik}

Marktlisten, Poolstände und Auszahlungsberechtigungen werden nicht in einer
Datenbank verwaltet, sondern aus der Chain gelesen. Ein klassischer Backend-Dienst
könnte höchstens als Cache, Indexer oder für Metadaten dienen; im vorliegenden
Repository ist bewusst darauf verzichtet, um die DApp-Eigenschaft nicht zu
relativieren. Die Ausnahme bildet die semantische Rolle der Admin-Adresse: Sie ist
kein Server, aber ein zentrales Vertrauenselement auf Protokollebene, weil sie
die Auflösungsentscheidung trifft (vgl. Kapitel~\ref{chap:sicherheit}).

\section{Netzwerke, Konfiguration und Betriebsmodi}

\paragraph{Lokale Entwicklung.}
Für den Alltag im Kurs eignet sich ein Hardhat-Node mit Chain-ID \texttt{31337}
und RPC-Endpunkt \texttt{http://127.0.0.1:8545}. MetaMask muss dieses Netzwerk
als benutzerdefinierte Chain eingetragen haben, damit Transaktionen nicht an das
falsche Netzwerk gesendet werden.

\paragraph{Öffentliches Testnet.}
Das Projekt unterstützt Sepolia-Deployments über Umgebungsvariablen wie
\texttt{SEPOLIA\_RPC\_URL} und \texttt{SEPOLIA\_PRIVATE\_KEY}. Private Schlüssel
dürfen niemals im Versionskontrolsystem, in CI-Logs oder in Container-Layern
landen; stattdessen sind Geheimnisse über die jeweilige Laufzeitumgebung zu
injizieren.

\paragraph{Konsistenz zwischen Frontend und Chain.}
Die Adressen in \path{frontend/app/contracts/addresses.json} müssen exakt zu der
Chain passen, gegen die MetaMask verbunden ist. Ein häufiger Fehler im Praktikum
ist die Kombination aus lokalem Node, aber noch ausstehendem
\texttt{export:frontend} nach einem frischen Deploy, oder ein Sepolia-Frontend mit
Hardhat-Adressen.
